% META: {"id":"1SPE-GEOREP-EX-030","chapitre":"1SPE-GEOMETRIE-REPEREE","type_objet":"exercice","capacites":["1SPE-GEOMETRIE-REPEREE-C3"],"capacites_codes":["C3"],"methodes":["M3"],"parcours":3,"competences":["calculer","raisonner"],"duree_min":15,"mode_creation":"ex_nihilo","sources_inspiration":[],"fichier_tex":"chapitres/1SPE-GEOMETRIE-REPEREE/exercices/1SPE-GEOREP-EX-030.tex","corrige_tex":"chapitres/1SPE-GEOMETRIE-REPEREE/corriges/1SPE-GEOREP-CO-030.tex","status":"approved","origin":"generated","status_history":["generated","audited","accepted","approved"],"release_acceptance":"RELEASE_OWNER_BATCH_ACCEPTANCE","acceptance_closure_digest":"sha256:9b3ccf9a81c5520fbb7e03b7b2d3e7bf2057a02834b6d908bf3be5f49f3b553f","acceptance_receipt_digest":"sha256:4fad86f0282e0fa4e0419cca1ad6379ae7af5a280c04a4a90880f90a1c044242"}
% BEGIN-VERIFY
% from sympy import *
% x, y = symbols('x y')
% # Montrer que M est sur le cercle de diametre [AB] ssi MA.MB = 0
% # A(0,0), B(6,2), M(x,y)
% # MA.MB = x(x-6) + y(y-2) = x^2-6x+y^2-2y
% # Cercle de diametre [AB]: centre I(3,1), r = sqrt(9+1)=sqrt(10)
% # (x-3)^2+(y-1)^2=10 => x^2-6x+9+y^2-2y+1=10 => x^2+y^2-6x-2y=0
% eq1 = x**2 - 6*x + y**2 - 2*y
% eq2 = (x-3)**2 + (y-1)**2 - 10
% assert expand(eq1 - eq2) == 0
% END-VERIFY
\begin{exercice}{1SPE-GEOREP-EX-030}{3}{15}
Soit $A(0 ; 0)$ et $B(6 ; 2)$. Montrer que l'ensemble des points $M(x ; y)$ tels que $\vec{MA} \cdot \vec{MB} = 0$ est un cercle dont on précisera le centre et le rayon.
\end{exercice}
